\section{The question and the proof mechanism}

A permutation matrix has one nonzero entry in every row and column. Its entries are therefore strongly dependent, even though each entry has mean $1/N$. This dependence has a useful consequence. Products of suitably separated centered entries are far smaller in expectation than one would predict by taking absolute values first.

For a uniform permutation matrix $M$ of size $N$, two distinct diagonal cells satisfy
\begin{equation}\label{eq:two}
\E\bigl[(M_{11}-N^{-1})(M_{22}-N^{-1})\bigr]
=\frac{1}{N(N-1)}-\frac{1}{N^2}
=\frac{1}{N^2(N-1)}.
\end{equation}
Two cells in the same row instead have centered product expectation $-N^{-2}$. The pattern of requested cells matters. The smaller scale in \eqref{eq:two} is cancellation, not independence.

The mixed-product estimate in Bordenave's Proposition 28 \cite{Bordenave} exploits this distinction. It also appears in Bordenave and Collins \cite[Proposition 22]{BC}. Brito, Dumitriu and Harris import it as Lemma 7 in their preprint, numbered Lemma 5.6 in the published article \cite{BDH}. We give a proof built from three operations:
\begin{enumerate}
\item Keep the isolated centered singletons together.
\item Condition only on a compatible set of prescribed matches. The remaining matrix is exactly a smaller uniform permutation matrix.
\item Use its row sum to obtain a positive recurrence, and only then expand the remaining factors.
\end{enumerate}
This order prevents the triangle inequality from destroying the cancellation we want to estimate.

For the diagonal product
\[
a_k(N)=\E\prod_{i=1}^k(M_{ii}-N^{-1}),
\]
Collins and Nagatsu establish the base recurrence, the parity-dependent leading terms, and the normalized error recurrences \cite[Proposition 3.3]{CN}. These are prior results. We reproduce their recurrence from a row-sum identity and use its positivity in the reverse direction, obtaining a lower bound that proves the cubic scale is necessary. A matching upper bound proves sufficiency. A continuous comparison then describes the transition itself.

\subsection{Main results at a glance}

Write $R_k=a_k/L_k$, where
\begin{equation}\label{eq:leading}
L_{2m}=(2m-1)!!N^{-3m},\qquad
L_{2m+1}=\frac{4m}{3}(2m+1)!!N^{-3m-2},\qquad m\ge1.
\end{equation}
The notation $(2m-1)!!$ means the product of the positive odd integers up to $2m-1$. The leading term for odd $k=1$ is zero, so $R_1$ is not defined.

The main conclusions are as follows.
\begin{itemize}
\item For a mixed product involving $a\le n$ distinct cells, with $h$ isolated centered singletons and $b$ nonisolated centered singletons,
\[
|\E P|\le2^b\frac{n^{-h}\mathcal P_h(n,a-h)}{(n)_a}.
\]
Here $\mathcal P_h$ is the explicit positive polynomial in \eqref{eq:depletedpoly}. When $b=0$, this is the exact maximum for fixed $n,a,h$, attained by compatible centered and raw cells in one permutation. If $b\ge1$, the stronger bound is
\[
|\E P|\le2^{b-1}\frac{n^{-h-1}\mathcal P_h(n,a-h-1)}{(n)_{a-1}}.
\]
For $b=1$ this is again an exact maximum, now attained by a pattern with one conflicting centered singleton. Theorem~\ref{thm:oneconflict} proves this stronger bound and its attainment when $b=1$. Scalar consequences include
\[
|\E P|\le\sqrt e\,2^b n^{-a}
\left(\frac{\sqrt{a-1}}{\sqrt n-\sqrt{a-1}}\right)^h,
\qquad 2\le a\le\sqrt n.
\]
Theorem~\ref{thm:mixed} allows arbitrary multiplicities and independent colors. Corollary~\ref{cor:mixedscalar} sharpens the singleton factor when $1\le h<a$ and proves optimality of the universal prefactor $\sqrt e$.
\item For every $2\le k\le N$,
\[
R_k-1\ge\frac{2k^3}{27N}.
\]
If $2k\le N$, then $R_k\le\exp(3k^3/N)$. Hence $R_k\to1$ exactly when $k^3/N\to0$, and $R_k=O(1)$ exactly when $k^3/N=O(1)$.
\item If $k^3/N\to c<\infty$, the even and odd ratios tend respectively to
\[
\cosh\left(\frac23\sqrt c\right),\qquad
\shc\left(\frac23\sqrt c\right),
\]
where $\shc(z)=\sinh(z)/z$ for $z\ne0$ and $\shc(0)=1$. Theorem~\ref{thm:critical} gives a uniform finite-size error.
\item With the exact parity-dependent $H_k$ in \eqref{eq:hyperdef}, the corrected formula is
\[
R_k\sim H_k\exp\{NG(\sqrt{k/N})\},\qquad
G(t)=\sum_{j\ge4}\frac{2t^j}{j(j-2)},
\]
uniformly throughout every sublinear upper range. Its first terms give sharp relative-accuracy thresholds $N^{1/2}$, $N^{3/5}$ and $N^{2/3}$. Theorem~\ref{thm:global} gives the approximation $\mathcal A_{N,k}$ at every order by keeping the denominator exact. Its optimal worst-case relative error is $C_*N^{-1/3}+o(N^{-1/3})$, where $C_*\approx0.39549627$. The full parity profile locates the largest error at odd orders near $1.1654N^{1/3}$.
\item Put $W_{N,k}=R_k/[e^{k^2/(4N)}H_k]$. For every fixed $C>0$, the exact worst-case error constant is
\[
\lim_{N\to\infty}N^{1/4}
\sup_{2\le k\le C\sqrt N}|W_{N,k}-1|
=\frac{\sqrt C}{2}+\frac{2}{15}C^{5/2}.
\]
Section~\ref{sec:rates} gives the signed parity expansion and an explicit finite bound.
\end{itemize}

The mixed estimate supplies the requested probabilistic input to the graph arguments. The moment asymptotics answer a separate quantitative question about that input. They do not by themselves improve the final spectral-gap theorem.
